\section{Eating your way to nondual awareness}

While Cologero is working on his little project and I am busily translating Sintesi, I'll toss out this short post regarding the Aghori\footnote{\url{http://www.ceveni.com/2009/01/aghori-sadhus-eating-human-flesh-video.html}} practice of cannibalism. Apparently, in order to reach nondual consciousness, the Agori practice the consumption and manipulation of semen, menstrual blood, faeces, and urine, all that we conventionally consider disgusting and filthy. That is what a fellow wrote about them; furthermore, he defended the practices since by transcending the categories of \emph{pleasant} and \emph{repellent}, the initiate is able to achieve nondual consciousness. Of course, the key term here is ``conventional'' ... are these categories merely a social construct or do they reflect a cosmic order? It should be obvious that we can then treat everything as a social construct and achieve a higher consciousness by flouting every category. What follows is a quick draft I wrote to participate in the discussion. I am interested in any feedback.

\begin{quotationx}
It is true that a man must overcome all those resentments, sentiments, opinions, etc, in him, that are arbitrary, instinctual, or chaotic. Nevertheless, that is not equivalent to accepting anything in him that represents formlessness or undifferentiatedness. Food requires cultivation and preparation, so it is part of cosmic order. Excrement, on the other hand, represents chaos and formlessness, so to equate the two leads to nothing. What is the transcendent unifying principle that overcomes that duality?

If cannibalism leads to anything, why is Jeffrey Dahmer not revered as a guru? If sex and intoxicants were so efficacious, then every inner city in the USA would be the destination of spiritual pilgrims. No, there needs to be something more to it.

If we follow the logic of the aghoris, then why would we distinguish between courage and cowardice? Is that an arbitrary duality? How about loyalty and treason? There is a reason that Arya means noble. The Western way means following the higher path, the more intelligent, the more artistic, the more difficult. Nondual awareness in the West means the golden mean. Courage is the unifying principle transcending recklessness and fear. Thus nondualism is transcendent and requires wisdom. There is neither transcendence nor wisdom required to be an Aghori ... it is not the path of the West.
\end{quotationx}



\flrightit{Posted on 2011-10-14 by Aeneas}

\begin{center}* * *\end{center}

\begin{footnotesize}
\begin{sffamily}
\texttt{Boreas on 2011-10-14 at 06:15 said:}

I find the Aghori LHP to be a very interesting phenomenon, yet I do agree that it is not suitable for most people, and even less for most westerners, and not the least because on a complete lack of a suitable social milieu and atmosphere. I have understood the Aghori path as a way for some rare people ``to experience God in horrendous forms'' and see everything as a manifestation of divinity -- even the most revolting aspects of manifestation.

It is very easy for a westerner (I'm not referring to you, Aeneas) to mis-understand the eastern tantra and left-hand path for a promiscuous and chaotic life-style with ``sex, drugs and rock n' roll'', while the truth is completely otherwise. If one needs to break taboos and social mores as an antinomian action, it should be done as a means to an end and not as an end in itself. Again, the intention behind any action counts the most. Most of the western LHP's mistake the higher and lower antinomianism and simply corrupt themselves and their possible followers morally, ethically, physically, psychically and socially. Yet, even if usually mis-understood and mis-applied, I'm not surprised that this kind of a life-style or antinomianism itself attracts those who have lived among anglo-saxon puritanism and victorian moralism. It is a search for a lost wholeness (I'm not referring to eating or drinking excrement, aha).

If I remember correctly, many of the Aghoris change later into righteous RHP brahmins if they pass thrgough the long and hard training period in the Aghori lifestyle. Many of them possibly dissolve into formlessness and chaos, but that seems to be the very challenge that Shiva offers them.

There's a good document about the Aghoris in the YouTube. Here's a link to the first part of the video:

\url{https://www.youtube.com/watch?v=W0bGrvKVxac}


\hfill

\texttt{Boreas on 2011-10-14 at 06:32 said:}

(Maybe someone could change that video-link into a more user-friendly mode? Thanks.)

The Aghori's remind me about the categories that Evola put forward in his Yoga of Power for those suitable to follow the LHP tantric path: the kaulas and viryas. Or at least these should be the requirements for anyone following the LHP.

\hfill

\texttt{HOO on 2011-10-14 at 07:55 said:}

Did the author of this post not even read relevant material such as ``The Yoga of Power'' before writing it? ``There is neither transcendence nor wisdom required to be an Aghori''. I promise you: you do not know that.

\hfill

\texttt{Will on 2011-10-14 at 10:07 said:}

Robert Svoboda's book Aghora: At the Left Hand of God is a very interesting account of the author's discipleship with an Aghori.

I agree with Boreas' comment. The majority of westerners who feel attracted to this sort of path will have all the wrong reasons for it. ``I take drugs and play music as my path to transcendence and stuff. I'm, like, spiritual.''

The basic distinction between the right and left hand paths are between renunciation and transformation. The right hand path renounces what is impure, while the left hand path alchemically transforms impurity into purity. I'm sure that the West has plenty examples of people who flouted conventional morality who were nonetheless spiritually adept. Dante's vision of Paradise is populated with some individuals who are rather questionable if viewed solely through the lens of morality.

In Buddhism, all tantric practices have the Mahayana ethical view (the desire to benefit all beings) as their foundation. In other words, you have to cultivate ethics and morality before you can transcend them. You have to be able to operate from the source of morality (which is Buddha nature or non-dual awareness) before you can start breaking the rules.

\hfill

\texttt{HOO on 2011-10-14 at 10:28 said:}

``I take drugs and play music as my path to transcendence and stuff. I'm, like, spiritual.'' is about as bad as ``I abstain from all socially-misunderstood drugs, and read books as my path to transcendence and stuff. I'm, like, spiritual.''

Some people who spend their whole life devotionally have no spiritual experience, while others who hardly make any effort at all, are powerful conductors of the spirit.

This is probably the harshest things for wannabes to deal with. And to have a place for their frustrations they form social circles in which they feel they have created a platform to judge their greaters from (``because this Traditionalist has this opinion too it must be right'').

\hfill

\texttt{Graham on 2011-10-14 at 12:58 said:}

Does anyone have a copy of the Pilgrim's Regress by CS Lewis? There's a part where the Pilgrim rebels against an evil giant who mocks his prisoners for preferring cow's milk to cow's manure. I can't remember how the argument went, exactly, but it reminds me of Aeneas' here.

\hfill

\texttt{invisible defender on 2011-10-14 at 22:05 said:}

``I take drugs and play music as my path to transcendence and stuff. I'm, like, spiritual.''

i agree with HOO, academic types are most likely to cynically dismiss this tendency. but setting aside the confusions and eccentricities for a moment, the popularity of such lifestyles can tell us things about our changing landscape.

sex, drugs and music are probably the three best ways of achieving visceral intensities capable of fracturing the solidification of residual materialism that still enclose people. for some this will be a way down, for a few it will also provide the fuel for an ascent. to see such habits take hold amongst ever greater numbers of young people reinforces what we have been told about the coming change in the world (for the worse, but change nevertheless).

i doubt whether anyone in the west can get anywhere today, in an operative sense, without first getting sucked in to the dirt of the psychic in some way, and fighting there to find the hidden energy that burns the illusion in some measure. but beware, a frail fight could see you routed into the reserve ranks with burlap garments, purple-lens glasses and dreadlocks! lolzolzlolz ;-)

\hfill

\texttt{Aeneas on 2011-10-14 at 23:37 said:}

It is true, HOO, I do not know any Aghori. Yet you read the letter of the post, but not the spirit, since you addressed nothing of substance. In context, a better way of putting it would be: ``There is neither transcendence nor wisdom required to take on the practices of an Aghori''. For that, I gave the example of Jeffrey Dahmer (for international readers, he was a homosexual serial killer who ate the bodies of his victims), and also the ghetto culture in the USA. So I concluded that ``No, there needs to be something more to it.'' It is embarassing to have to explain it, as if to someone who missed a joke. So what exactly is it above and beyond cannibalism? That is a legitimate question to ask, nay, it is the one question necessary to ask. And to that I got no response, other than to go read a book ... from someone who mocks readers.

So, HOO, perhaps you can share your own experiences: when you perform fallatio on your friends, does it lead to enlightenment? Assuming, of course, the thought of it repulses you.

\hfill

\texttt{Aeneas on 2011-10-14 at 23:45 said:}

The battle cry of the 68er was ``sex, drugs, and rock n roll''. The results of more than 40 years of that is all too obvious, all least to those who have eyes to see and ears to hear. As for getting sucked into the psychic dirt, that takes absolutely no effort, you just need to wake up in the morning. Here we follow the Aryan path of the greater battle: to overcome all that psychic dirt within us.

Of course for the mass man, visceral experience is all he can strive for. So he praises sex, drugs, and music because that is the democratization of all experience, since it is accessible to all and sundry. Go for it, dude, but the results are certain.

The comments to this post are so pathetic, they make me want to puke. Wow, I feel much more enlightened! in the operative sense, of course. Time for another jello shot.

\hfill

\texttt{Utlage on 2011-10-15 at 00:09 said:}

Okay, let's not lose our minds here. The Aghori represent an extreme degeneration of even a lower-grade spirituality.

True, yes, yes, there are ``Enlightened Lunatics'', the theme is well-known... Within Hinduism there is no absence of orthodox Hindus basically excommunicating and condemning the Aghori for their intellectual heterodoxies, leading to their transgressive self-bestiality... But let's call a spade a spade here: these people are not necessarily ``enlightened'' simply because outrageously antinomian--psychiatrically, most likely they are sufferers of mental disorders; or if not that, as aforementioned, intellectual errors condemned by orthodox Hindus allowing such subhumanization to arise... 

\hfill

\texttt{EXIT on 2011-10-15 at 10:01 said:} 

\emph{Not until after 5, dude, when we have company. But thanks for letting us know what you do with your solitary time. --- admin} You guys at gornahoor don't drink, have sex, and listen to music? No wonder you're so uptight.

\hfill

\texttt{HOO on 2011-10-15 at 11:18 said:} 

Evola: ``essentially, one must be radically anti-bourgeois and oppose all bourgeois compromises and conformity. Years ago, Ernst Jünger wrote something that he could never agree with today (for he, too, the decorated combatant, has reformed and fallen in line). Jünger wrote: `Better a delinquent than a bourgeois.'\,'' Evola, `Incontro con Julius Evola', L'Italiano 11 (November 1970, pp. 812-20). -- invisible defender: '' but beware, a frail fight could see you routed into the reserve ranks with burlap garments, purple-lens glasses and dreadlocks!''

I doubt that's the case for anybody who isn't so predisposed. Predisposition is of utmost momentousness.

One of the few things I truly know is that I'll never be a part of the horde you speak of. I couldn't if by my life I wanted or tried. The repel against me like bugs from fire, even if I try entryism. If I succeed for a moment, they eventually sense there is something very, very different about me. My splendor is their horror.

You really seem to emphasize this fear ``of succumbing to the horde'' when you speak to me. I've never feared it because I know it's all in gods hands; ``once saved always saved''. Maybe you fear it and project upon me for the reason that, unlike me, you actually are in danger of succumbing? : )

``They whom God hath accepted in His Beloved, effectually called and sanctified by his Spirit, can neither totally nor finally fall away from the state of grace; but shall certainly persevere therein to the end, and be eternally saved.'' Westminster Confession of Faith (chap. 17, sec. 1) 

Nothing can kill my interest in ``the primal source of all, the infinite, mysterious, and immutable love of God.''

Those of us who develop the capacity to accept that which is unacceptable inside of us, those of us who develop the capacity to fully hold in our hearts, our bodies, and our minds that which we have always both consciously and unconsciously rejected and run from, will be those who will break through into the depths of our true being.

The other comments here are littered by signifiers/signals of weakness, frailty, fear, ideology, dogma... they are littered with the signs of a lot that lacks experience of the world of god. Aeneas: ``The comments to this post are so pathetic, they make me want to puke.``

Aeneas: `` So, HOO, perhaps you can share your own experiences: when you perform fallatio on your friends, does it lead to enlightenment?{\kern0pt}``

Aeneas: ``The comments to this post are so pathetic, they make me want to puke.``

Further, the comments here, such as made by Utlage, are very ineffectual in dispelling everything else written about Tantrism by Evola and other authors. You lot, you socialist (for ye meddle with the affairs of others as if out of instinct) brats, are just constructing your platform of mediocrity of which the building material is opinions. You are building on sand. Evola: ``What I regard as the most important thing in a person is his `existential reality'\,''!

Indeed.

\hfill

\texttt{Boreas on 2011-10-15 at 16:52 said:}

I'm sorry if you find (my) comments pathetic Aeneas, but you asked for any kind of feedback. My own knowledge of Aghoris is limited to generalities and to that document I linked. I know a friend who has read Svoboda's book about them, and his opinion of the book was that there were some glimpses of real understading of the LHP and ascetic practices (detachment etc.), but in general the book was quite low-grade and somewhat sensational -- which isn't really a surprise.

``So what exactly is it above and beyond cannibalism? That is a legitimate question to ask, nay, it is the one question necessary to ask.''

That's true. I really can't see the practices of the Aghori sect being limited or being reducible to cannibalism. It seems to me that most of their practices are basic practices of the Shaivite-kind mixed with elements of the tantric LHP and Kali worship. I think another legitimate question in general to ask about eating human flesh from cremation pyres is that what differentiates human flesh from any other meat?

Psychic residues of the particular human stock might be one answer, and perhaps this could lead into chaotic psychic dis-orders very easily, but looked from another angle, in rare cases and in ``initiatic consciousness'' that meets certain requirements, it could also lead into the transformation of these psychic elements into magical powers, or at least into a heightened vitality, in the same manner than in tantric practices (sex \& intoxicants). Lets remember that the brain-eating of dead enemies was one of the ``demonic-initiatic'' practices among savage people that was believed to transfer the ``mana'' of the dead enemy to the one who eated them; maybe in another context the demonic element could be transformed into something higher?

Physically speaking, the obtained results could also be two-fold: either some sort of disease in the likeness of a ``mad cow disease'' (which is a result of consuming nutrition from the same species, ``too close''), or a heightened physical immunity to diseases and viruses.

(Again, these are speculations deriving from my background as an nutritionist who has studied esoteric doctrines. I hope they are worth more than nothing.)

I think there are differences among Aghoris in the cannibalism question. Most certainly there are also those who do it simply because they are corrupted men in search of black psycho-magical powers etc., but we should not exclude other possibilities and intentions out of hand. In the document I linked the particular Aghori (can't remember his name right now) simply expressed that the nutrition he eated could also be human flesh (or feces), but did not especially strive for these as foods. He simply used, ate and consumed everything that was around in the cremation grounds, including flesh of dead corpses. The urge he explained about that occurred in him in this instance may be an expression of some sort of a primitive atavistic bloodlust, but it can be something other also, or at least transferable to something other than this. Can't really tell.

``So, HOO, perhaps you can share your own experiences: when you perform fallatio on your friends, does it lead to enlightenment? Assuming, of course, the thought of it repulses you.''

I am not HOO, but I'll answer to this in general the following way: I think you are lumpng different things / categories together here. The Aghoris do not eat human flesh or feces to humiliate themselves willfully or to show off. Neither they do it for enlightenment, at least not simply for this goal in mind. If I have understood anything at all about their motives, it could be expressed in the way of the ``ishmaelite'' expression ``nothing is sacred, everything is allowed'', or in some similar manner that describes this state of mind. Maybe it is a practice of ``assuming God-forms'', which would mean in this case the assumption of the third face of Shiva that consumes everything and makes no difference between the body of a whore, king, ascetic or dog.

\hfill

\texttt{Aeneas on 2011-10-15 at 22:05 said:}

Sorry, Boreas, for the misunderstanding ... I was responding to Invisible Defender.

Why do we not eat the dead? Man ist was er isst.

If we don't honour the dead, then we won't honour their tradition. After all that has been written here about ancestors in the Ancient City, this train of thought is diappointing. The play Antigone demonstrates the Western way to treat the dead, not the Aghoris. The mark of the intelligent man is the ability to discriminate. So I repeat my previous assertion: there is no intelligence in nihilism, since the effort to make distinctions is lacking. I'll ask the question again: why do we praise the hero and revile the coward? Is that a difference unknown to Shiva? Why does Arjuna do his duty if it is not a sacred vow and everything is allowed?

No one is thinking this through.

\hfill

\texttt{Boreas on 2011-10-16 at 05:50 said:}

What you're saying is absolutely true Aeneas, but no one has proposed that we should all act like Aghoris, least of all the Aghoris themselves, or that it could be applied by westerners without serious problems. Their path -- whether it is a legitimate one or not -- is for a very restricted number of people. Without better knowledge I suspect that even most of the Shaivites look upon them as heretics.

As you mentioned Arjuna, wouldn't the actions dictated by Khrisna to him be in another context the very epitome of dis-honour? ``Kill your relatives and friends'' etc. In another context they become part of Arjuna's dharma, the violation which would bring him the most severe act of penance. In this Aghori case we are talking antinomian a-dharma, that's true, but this is after all the very credo of the tantric LHP. Their very methods of liberation are such that they would bring disaster to a pashu or even for a ordinary initiate.

In another cultural context it could very well be an act of honouring the dead to use or even consume their material remains as throughly as possible. It is after all only their body, their material counterpart, and the consumption of their flesh could be linked to the continuation of their life-force in the next tribe leader etc. I agree this sounds rather primitive and savage, but I'm only trying to understand in what context the consuming of human meat could be linked to some sort of sacrality or initiatic practices. Severe restrictions and discriminations are certainly needed.

\hfill

\texttt{Will on 2011-10-16 at 11:37 said:}

The distinction between doing one's duty and forsaking it is important, but according to tradition, there are those few who genuinely transcend it. To forsake one's duty while still believing in duty and still having a karmic bond to the realm in which one has a duty is to reap the negative consequences of that act of neglect. But it is possible -- according to doctrine, at least -- to transcend the very notion of duty, such that one remains stainless no matter what actions he undertakes, because there is no longer a `person' who identifies with the `doer' of the actions.

Krishna told Arjuna to do his duty and fight, but he also told him not to identify with any actions, so as not to reap the karma of killing and fighting. As Robert Svobodha's Aghori teacher told him, ``The moment you act to protect what you feel to be your self-interest, your ego, karma adheres to you like mud.'' The point is not that war and killing are good per se, but that all phenomena are primordially pure, even war and killing.

You can see how the notion of transcending one's duty is dangerous, and how it has wreaked havoc in the West. When the 60s generation became interested in Buddhism and Hinduism, they were attracted to this idea of transcending their duty -- their duty to family, society, and the world. But without genuine egolessness, there is no transcending duty -- there is only forsaking duty. I personally think that people in the West need to put down the high falutin Eastern metaphysics for a while and go back to basics with Marcus Aurelius and the four cardinal virtues.

\hfill

\texttt{Kaulaphon on 2011-10-16 at 17:35 said:}

Some feedback:

As a practicioner of vamachara tantra, I actually believe that the 'right-hand path' of achievement is far more dangerous in the end. The sudden path have a paradoxical safeguard in that the failures soon becomes demented and crazed -- while those on the long trek towards Mt. Meru have every opportunity to weave themselves into intricate webs of abstractions whose threads tend to disentangle the individual from the direct experience of spiritual potency, replacing it instead with abstractions, mind games, dogma, subtle safety and complacent comfort where the available bodily, emotional and mental energy are invested away into what in the end often translates as nothing more than spiritual ponzi schemes.

So, on a practical level, what is the meaning of 'cosmic order'? First, one must understand that the tantrika perceives the experience of dual reality not as some illusory trickery to escape from, but as a cosmic web of eternal power. The principle behind this understanding is the realization of 'svatantrya' (self-dependency) -- the ability to influence any agent while oneself remains untainted by any such influence. This is the deepest nature of shakti -- of power -- and of left-handed attainment.

Now, mostly, shakti is a mean bitch. She binds experience, she dominates, she subjugates. This is why most of mankind are regarded as 'pashu' -- as beasts of burden -- as their whole experience is to dance as tamed animals to the sensory and instinctual yoke life pulls them along with. One can therefore approach power from three angles: to submit, to escape, to overcome.

The Left Hand Path is all about overcoming. Not to avoid experience, but to confront it, and to master it. Not to submit to experience, but to transmute it.

What is 'order' then -- understood from an alchemical and traditional perspective? Order is form, it is an effort of imposing structure, of incarnating unbound will in bound circumstances. This IS the magician's task: to perceive any experience in the Universe as a tranformative symbol, as some 'will' given appearance, as a symbol that may be manipulated and interacted with, dissolved and synthesised, and submitted to the will of the magician. The magician is therefore, to paraphrase Evola, a god, a builder of worlds -- a cakravartin -- ruler of worlds.

And the world he is building is his own, the cosmic order of the Self. In traditional societies, all other order and hierarchy stems from this original imposing of order by the magical elite. Or so speaks the mythos. The task itself is reflected in societal functions: the magician issues the will, the warrior honors, defends, expands and uphelds it, and the producer and the farmer and the craftsman shape the base world according to it. Thus is limitless will incarnated, made manifest as limited form. And it matters not that dual form is in the final analysis ephemeral (or chaotic, or 'evil' etc), for the non-dual will that dominate form is not.

Breaking a taboo to overcome socially conditioned morality -- eating human flesh etc -- is just one facet of the tantrik principle of viparit-karani ('opposite doing'). By truly overcoming the taboo, energy (power) is liberated -- and the practioner is no longer bound by its chains. But on a deeper level, there are alchemical considerations as well. Therefore we have whispered rumours of certain tantric sects amongst the tibetans that engage in the purposeful eating of spiritual accomplised humans, some times by force. Not for the taboo, but for the attainment of certain alchemical properties. Be that as it may, the underlying point is tied to the biospiritual practice of advanced tantrik yoga: the utilization of energetic forms otherwise shunned, to manipulate the subtle bodies of the practicioner in a higly efficient manner in order to achieve illuminated power in both flesh and spirit. Electricity serves an analogy: a strong and otherwise deadly current is trapped, controlled and tapped by those who have the knowledge and the tools to deal with such risky business.

\hfill

\texttt{Boreas on 2011-10-17 at 03:56 said:}

Thanks for your insights Kaulaphon, they were illuminating to say the least. I didn't know about the alchemical practices of the tibetan tantrics you mentioned, that was basically what I was referring to in my previous imaginary example of the possible sacral-magical eating of the bodies of the deceased.

\hfill

\texttt{Aeneas on 2011-11-04 at 12:49 said:}

Apparently the New York City school system is providing Aghori training for its students. However, they call it \textit{scat play}\footnote{\url{http://www.lifesitenews.com/news/recommended-new-york-sex-ed-site-offers-middleschoolers-tips-on-eating-fece}}.

\hfill

\texttt{poopalate on 2011-11-05 at 08:50 said:}

Yep, the Golden Age is riiiiight around the corner.

\hfill

\texttt{Boreas on 2011-11-16 at 06:43 said:}

Aeneas wrote: ``Apparently the New York City school system is providing Aghori training for its students.''

``When the Tao was lost... ''

Thank Lord `The Age of Universal Bullshit' will come to an end eventually.

\hfill

\texttt{Aldebaran on 2012-02-05 at 09:48 said:}

Aeneas:

``No one is thinking this through''

Perhaps not, but you aren't thinking, at all. Your response to this subject is purely emotional. You ask for comments, and then become personally insulting when the content isn't to your liking.

In addition, your form of ``argument'' mirrors that of many Traditionalists: The tactic of merely looking down your nose through your monocle with contempt at anything you find repugnant (or don't understand; often one and the same), and hoping that others won't notice the complete absence of logic, argument, or evidence. I've noticed that this is a favorite tactic of Evola's, in particular, so you clearly come by it honestly.

\hfill

\texttt{Cologero on 2012-02-20 at 20:26 said:}

I don't understand the hostility, Aldebaran. If you have a logical argument or new evidence to shed light on this issue, please share it with us.

\hfill

\texttt{Perennial on 2012-02-21 at 00:44 said:}

Ah, nothing like the old ad hominem.


\hfill

\end{sffamily}
\end{footnotesize}

